\subsection*{А.1. Промпт для перевода полей \texttt{documents\_sentences}}

Ниже приведён системный и пользовательский промпты, использованные в Qwen2.5-72B-Instruct для перевода полей \texttt{documents\_sentences} (переведены на финальной версии конвейера, см. раздел~\ref{sec:translation_pipeline}).

\begin{verbatim}
SYSTEM:
You translate from English to Russian. The input is a JSON list
of sentences like [["0a","..."],["0b","..."], ...]. Return the SAME
JSON list, where Russian translations keep the same positions.
Rules:
- DO NOT change the number of list items.
- DO NOT change the keys (the first element in each pair).
- DO NOT add Markdown, numbering, comments, or explanations.
- Preserve proper names, numbers, measurement units, and URLs.
- Translate fluently without changing the meaning.

USER:
<json-array-of-sentences>
\end{verbatim}

\subsection*{А.2. Промпт для LLM-разметки (Qwen2.5-72B)}

\begin{verbatim}
SYSTEM:
You annotate the quality of RAG systems. Determine:
1) which document sentences are relevant to the question;
2) which of them are actually used in the answer;
3) whether the full answer is supported by the context (adherence, true/false).

USER:
Question: {question}
Documents (JSON, list of sentence lists with keys):
{documents_sentences}
Answer: {response}

Return JSON:
{
  "all_relevant_sentence_keys":  [<keys>],
  "all_utilized_sentence_keys":  [<keys>],
  "adherence_score":             true | false
}
\end{verbatim}

\subsection*{А.3. Статистика ошибок LLM-оценщика по поддатасетам}

Таблица~\ref{tab:llm_errors} - агрегаты ошибок разметки Qwen2.5-72B 
(overlap, miss, extra) на 7 тестовых поддатасетах RAGBench-RU. \emph{overlap} - 
доля ключей, совпавших с эталоном; \emph{miss} - среднее число пропущенных ключей 
на пример; \emph{extra} - среднее число лишних.

\begin{table}[H]
\centering
\small
\caption{Поддатасетные метрики LLM-оценщика Qwen2.5-72B}
\label{tab:llm_errors}
\begin{tabular}{lcccccc}
\toprule
поддатасет & rel.~overlap & rel.~miss & rel.~extra & util.~overlap & 
util.~miss & util.~extra \\
\midrule
hotpotqa   & 0.74 & 0.87 & 1.12 & 0.85 & 0.41 & 0.68 \\
hagrid     & 0.68 & 1.25 & 1.34 & 0.81 & 0.55 & 0.73 \\
expertqa   & 0.65 & 1.43 & 1.51 & 0.79 & 0.60 & 0.81 \\
covidqa    & 0.71 & 0.94 & 1.08 & 0.84 & 0.44 & 0.67 \\
delucionqa & 0.76 & 0.71 & 0.95 & 0.87 & 0.38 & 0.61 \\
cuad       & 0.62 & 1.68 & 1.82 & 0.76 & 0.72 & 0.94 \\
finqa      & 0.72 & 0.99 & 1.15 & 0.85 & 0.43 & 0.69 \\
\midrule
\textbf{среднее} & \textbf{0.70} & \textbf{1.12} & \textbf{1.28} & \textbf{0.83}
 & \textbf{0.50} & \textbf{0.73} \\
\bottomrule
\end{tabular}
\end{table}

Наибольшая доля избыточной генерации наблюдается в \texttt{cuad} (юридический домен)
и \texttt{expertqa} (научный домен), где LLM склонна добавлять
<<похоже релевантные>> предложения в разметку.

\subsection*{А.4. Полная сетка конфигураций}

Полная таблица всех $\sim$80 обученных конфигураций 
(базовая модель, длина, веса, язык, голова, F1 на уровне токенов по таргетам) 
доступна в файле \texttt{analysis\_output/tables/full\_grid.csv}. 
В настоящем тексте приведены только основные агрегаты; CSV распространяется вместе 
с кодом работы.

\subsection*{А.5. Пример конфига эксперимента}

\begin{lstlisting}[basicstyle=\ttfamily\scriptsize,breaklines=true]
# model_utils/config.py (fragment)
@dataclass
class DataConfig:
    question_key:                str = "question_ru"
    documents_sentences_key:     str = "documents_sentences_ru"
    response_key:                str = "response_ru"
    relevant_keys_field:         str = "all_relevant_sentence_keys"
    utilized_keys_field:         str = "all_utilized_sentence_keys"
    adherence_field:             str = "adherence_score"
    max_length:                  int = 2048

@dataclass
class TrainingConfig:
    num_epochs:            int   = 2
    train_batch_size:      int   = 64
    eval_batch_size:       int   = 32
    backbone_lr:           float = 2e-5
    head_lr:               float = 3e-4
    weight_decay:          float = 1e-3
    grad_clip:             float = 1.0
    warmup_ratio:          float = 0.1
    loss_weights:          Tuple[float,float,float] = (3.0, 3.0, 1.0)
    amp:                   bool  = True
\end{lstlisting}
